\section*{Capítulo \ref{ch:b2:integration} --- Integración}

\begin{solution}{exo:b2:integration:1}
$\int_0^1 \frac{\dd t}{\sqrt{t(1-t)}}$: cerca de $0$ es $\sim
t^{-1/2}$ ($\alpha = \frac12 < 1$: \omterm{def:b2:integration:improper}{converge}); cerca de $1$ es $\sim
(1-t)^{-1/2}$: \omterm{def:b2:integration:improper}{converge}. Convergente (su valor es $\pi$, mediante la
sustitución $t = \sin^2\theta$).

$\int_1^\infty \frac{\ln t}{t^2}$: $\frac{\ln t}{t^2} =
o(t^{-3/2})$: convergente (de valor $1$, por partes).

$\int_0^\infty \frac{\dd t}{1 + t^2\sin^2 t}$: divergente. Cerca de
$t = n\pi$, escribamos $t = n\pi + u$: $\sin^2 t = \sin^2 u \leq
u^2$, así que para $\abs u \leq \frac{1}{n}$ se tiene $1 + t^2\sin^2
t \leq 1 + (n\pi + 1)^2u^2 \leq C n^2 u^2 + 1$; por tanto
\[
\int_{n\pi - 1/n}^{n\pi + 1/n} \frac{\dd t}{1 + t^2\sin^2 t}
\geq \int_{-1/n}^{1/n} \frac{\dd u}{1 + Cn^2u^2}
= \frac{2\arctan\sqrt C}{\sqrt C}\cdot\frac{1}{n} ,
\]
término de una serie divergente de tipo armónico: sumando sobre $n$,
la primitiva no está acotada.
\end{solution}

\begin{solution}{exo:b2:integration:2}
Sustituyendo $u = \lambda t$:
\[
\int_0^\infty t^n \eu^{-\lambda t}\dd t
= \frac{1}{\lambda^{n+1}}\int_0^\infty u^n\eu^{-u}\dd u
= \frac{\Gamma(n+1)}{\lambda^{n+1}} = \frac{n!}{\lambda^{n+1}} .
\]
Con $t = \eu^{-u}$ ($\dd t = -\eu^{-u}\dd u$):
\[
\int_0^1 (\ln t)^n \dd t = \int_0^{\infty} (-u)^n \eu^{-u}\,\dd u
= (-1)^n\, n! .
\]
\end{solution}

\begin{solution}{exo:b2:integration:3}
Convergencia: el integrando es $\leq \frac{1}{1+t^2}$ cerca de
$\infty$ y está acotado cerca de $0$ (ambos factores están acotados
inferiormente lejos de $0$): \omterm{def:b2:series:def}{absolutamente} convergente para todo
$x$. Sustituyendo $t = \frac1u$ ($\dd t = -\frac{\dd u}{u^2}$):
\[
I(x) = \int_0^\infty \frac{1}{\bigl(1 +
\frac1{u^2}\bigr)\bigl(1 + u^{-x}\bigr)}\cdot\frac{\dd u}{u^2}
= \int_0^\infty \frac{u^x}{(1 + u^2)(1 + u^x)}\,\dd u .
\]
Sumando las dos expresiones de $I(x)$:
\[
2I(x) = \int_0^\infty \frac{1 + t^x}{(1+t^2)(1+t^x)}\dd t
= \int_0^\infty \frac{\dd t}{1 + t^2} = \frac{\pi}{2} :
\]
$I(x) = \frac\pi4$, independiente de $x$.
\end{solution}

\begin{solution}{exo:b2:integration:4}
Cerca de $0^+$, con $u = \abs{\ln t} \to \infty$. Si $\alpha < 1$:
convergencia sea cual sea $\beta$ (compárese con $t^{-\alpha'}$ para
$\alpha < \alpha' < 1$: el factor logarítmico queda vencido). Si
$\alpha > 1$: divergencia sea cual sea $\beta$ (compárese con
$t^{-\alpha''}$ con $1 < \alpha'' < \alpha$). Si $\alpha = 1$:
sustitúyase $t = \eu^{-u}$:
\[
\int_0^{1/2} \frac{\dd t}{t\,\abs{\ln t}^\beta}
= \int_{\ln 2}^{\infty} \frac{\dd u}{u^\beta},
\]
convergente si y solo si $\beta > 1$. En resumen: \omterm{def:b2:integration:improper}{converge} si y solo
si $\alpha < 1$, o bien si $\alpha = 1$ y $\beta > 1$; el espejo de
las series de Bertrand.
\end{solution}

\begin{solution}{exo:b2:integration:5}
\omterm{def:b2:metric:continuity}{Continuidad} sobre $\intco{0}{\infty}$: dominación
$\bigl|\frac{\eu^{-xt}}{1+t^2}\bigr| \leq \frac{1}{1+t^2}$,
integrable y uniforme en $x \geq 0$:
\cref{thm:b2:integration:continuity}.

$C^2$ sobre $\intoo{0}{\infty}$: para $x \geq a > 0$, las dos
primeras derivadas en $x$, $\frac{-t\,\eu^{-xt}}{1+t^2}$ y
$\frac{t^2\eu^{-xt}}{1+t^2}$, están dominadas por $t\,\eu^{-at}$ y
$\eu^{-at}$: dos aplicaciones del~\cref{thm:b2:integration:leibnizrule}. Entonces
\[
F''(x) + F(x) = \int_0^\infty \frac{t^2 + 1}{1 + t^2}\,\eu^{-xt}\dd
t = \int_0^\infty \eu^{-xt}\dd t = \frac1x .
\]
Límite: $0 \leq F(x) \leq \int_0^\infty \eu^{-xt}\dd t = \frac1x \to
0$.
\end{solution}

\begin{solution}{exo:b2:integration:6}
Sobre $\intcc{\varepsilon}{M}$, sustituyamos $u = at$ y $u = bt$ en
las dos mitades:
\[
\int_\varepsilon^M \frac{f(at) - f(bt)}{t}\dd t
= \int_{a\varepsilon}^{aM}\frac{f(u)}{u}\dd u
- \int_{b\varepsilon}^{bM}\frac{f(u)}{u}\dd u
= \int_{a\varepsilon}^{b\varepsilon} \frac{f(u)}{u}\dd u
- \int_{aM}^{bM} \frac{f(u)}{u}\dd u .
\]
Primera pieza: $f(u) = f(0) + o(1)$ cerca de $0$, y
$\int_{a\varepsilon}^{b\varepsilon} \frac{\dd u}{u} = \ln\frac ba$:
la pieza tiende a $f(0)\ln\frac ba$. Segunda pieza: $f(u) \to
f(\infty)$, mismo cálculo: tiende a $f(\infty)\ln\frac ba$. Por
tanto, la \omterm{def:b2:integration:improper}{integral impropia} \omterm{def:b2:integration:improper}{converge} a $\bigl(f(0) -
f(\infty)\bigr)\ln\frac ba$.

Con $f(t) = \eu^{-t}$ ($f(0) = 1$, $f(\infty) = 0$), $a = 1$ y $b =
2$:
\[
\int_0^\infty \frac{\eu^{-t} - \eu^{-2t}}{t}\dd t = \ln 2 .
\]
\end{solution}

\begin{solution}{exo:b2:integration:7}
Extendamos el integrando por $0$ más allá de $t = n$: $g_n(t) = (1 -
\frac tn)^n t^{x-1}\mathbf{1}_{t \leq n}$. Puntualmente, $g_n(t) \to
\eu^{-t}t^{x-1}$ (límite del interés compuesto, volumen del primer
año). Dominación: $\ln(1 - u) \leq -u$ da $(1 - \frac tn)^n \leq
\eu^{-t}$ sobre $\intcc{0}{n}$, de modo que $\abs{g_n(t)} \leq
\eu^{-t}t^{x-1} = \varphi(t)$, integrable. Convergencia dominada:
\[
\int_0^n \Bigl(1 - \frac tn\Bigr)^n t^{x-1}\dd t
\xrightarrow[n\to\infty]{} \int_0^\infty \eu^{-t}t^{x-1}\dd t
= \Gamma(x) .
\]
(Calcular el miembro izquierdo por partes reiteradas da la forma de
producto de Euler $\Gamma(x) = \lim
\frac{n!\,n^x}{x(x+1)\cdots(x+n)}$.)
\end{solution}

\begin{solution}{exo:b2:integration:8}
$H$ es derivable en $x$ (integrando de clase $C^1$ en $x$, con
derivada $-2x(1+t^2)\cdot\frac{\eu^{-x^2(1+t^2)}}{1+t^2} =
-2x\,\eu^{-x^2}\eu^{-x^2t^2}$, \omterm{def:b2:metric:continuity}{continua} y acotada sobre los
\omterm{def:b2:metric:compact}{compactos} de $x$, con dominación trivial sobre $t \in
\intcc{0}{1}$):
\[
H'(x) = -2x\,\eu^{-x^2}\int_0^1 \eu^{-x^2t^2}\,\dd t
\overset{u = xt}{=} -2\,\eu^{-x^2}\int_0^x \eu^{-u^2}\,\dd u
= -G'(x),
\]
ya que $G'(x) = 2\eu^{-x^2}\int_0^x \eu^{-t^2}\dd t$ (regla de la
cadena sobre el cuadrado y teorema fundamental del cálculo). Luego
$G + H$ es constante e igual a $G(0) + H(0) = 0 + \int_0^1 \frac{\dd
t}{1+t^2} = \frac\pi4$.

Cuando $x \to \infty$: $0 \leq H(x) \leq \eu^{-x^2}\int_0^1 \dd t
\to 0$, de modo que $G(x) \to \frac\pi4$:
\[
\int_0^\infty \eu^{-t^2}\dd t = \sqrt{\frac\pi4} =
\frac{\sqrt\pi}{2} .
\]
(En consecuencia, $\Gamma\bigl(\frac12\bigr) = 2\int_0^\infty
\eu^{-t^2}\dd t = \sqrt\pi$, por la sustitución $t = \sqrt u$.)
\end{solution}

\begin{solution}{exo:b2:integration:9}
Cauchy--Schwarz (volumen del primer año, válida sobre
$\intcc{\varepsilon}{M}$ y pasada al límite) aplicada a la
factorización $t^{\frac{x+y}{2}-1}\eu^{-t} =
\bigl(t^{x-1}\eu^{-t}\bigr)^{1/2} \bigl(t^{y-1}\eu^{-t}\bigr)^{1/2}$:
\[
\Gamma\Bigl(\frac{x+y}{2}\Bigr) \leq
\Gamma(x)^{1/2}\,\Gamma(y)^{1/2}
\quad\Longrightarrow\quad
\ln\Gamma\Bigl(\frac{x+y}{2}\Bigr) \leq \frac{\ln\Gamma(x) +
\ln\Gamma(y)}{2} :
\]
$\ln\Gamma$ es convexa en el punto medio y, siendo \omterm{def:b2:metric:continuity}{continua}
(\cref{thm:b2:integration:gammaprops}), es convexa
(\cref{exo:b2:realfun:8}). (La log-convexidad fija $\Gamma$ de
manera única entre las interpolaciones del factorial: el teorema de
Bohr--Mollerup, una perla del tercer año.)
\end{solution}

\begin{solution}{exo:b2:integration:10}
\begin{enumerate}
  \item Para $x \geq a > 0$: la derivada en $x$ del integrando es
        $-\sin t\,\eu^{-xt}$, dominada por $\eu^{-at}$; el~\cref{thm:b2:integration:leibnizrule} da $F'(x) =
        -\int_0^\infty \eu^{-xt}\sin t\,\dd t$. Dos integraciones
        por partes (o la exponencial compleja):
        \[
        \int_0^\infty \eu^{-xt}\sin t\,\dd t
        = \Im \int_0^\infty \eu^{(-x+\iu)t}\dd t
        = \Im\frac{1}{x - \iu} = \frac{1}{1 + x^2} .
        \]
  \item $\abs{F(x)} \leq \int_0^\infty \eu^{-xt}\dd t = \frac1x \to
        0$. Integrando $F' = -\frac{1}{1+x^2}$ de $x$ a $\infty$: $0
        - F(x) = -\bigl(\frac\pi2 - \arctan x\bigr)$, luego $F(x) =
        \frac\pi2 - \arctan x$.
  \item Haciendo $x \to 0^+$ con la \omterm{def:b2:metric:continuity}{continuidad} admitida: $F(0^+) =
        \frac\pi2$, y $F(0) = \int_0^\infty \frac{\sin t}{t}\dd t$
        (la integral semiconvergente de Dirichlet,
        \cref{ex:b2:integration:sint}): su valor es $\frac\pi2$.
\end{enumerate}
\end{solution}

\begin{solution}{exo:b2:integration:11}
Convergencia: cerca de $0$ el integrando se extiende
continuamente con el valor $1$ ($\sin t \sim t$); en el infinito es
$\leq t^{-2}$: convergencia absoluta. Sobre
$\intcc{\varepsilon}{M}$, integremos por partes con $u = \sin^2 t$ y
$v' = t^{-2}$:
\[
\int_\varepsilon^M \frac{\sin^2 t}{t^2}\dd t
= \Bigl[-\frac{\sin^2 t}{t}\Bigr]_\varepsilon^M
+ \int_\varepsilon^M \frac{2\sin t\cos t}{t}\dd t
= \Bigl[-\frac{\sin^2 t}{t}\Bigr]_\varepsilon^M
+ \int_{2\varepsilon}^{2M} \frac{\sin u}{u}\dd u
\]
($u = 2t$ en la última integral). El corchete tiende a $0$ en ambos
extremos ($\sin^2\varepsilon/\varepsilon \leq \varepsilon$; $\sin^2
M/M \leq 1/M$), y la última integral tiende a $\int_0^\infty
\frac{\sin u}{u}\dd u = \frac\pi2$ (\cref{exo:b2:integration:10}).
Por tanto
\[
\int_0^\infty \Bigl(\frac{\sin t}{t}\Bigr)^{\!2}\dd t
= \frac\pi2 .
\]
\end{solution}

\begin{solution}{exo:b2:integration:12}
\begin{enumerate}
  \item Por partes con $u = \frac{-1}{2t}$ y $v' = -2t\,\eu^{-t^2}$
        (de modo que $v = \eu^{-t^2}$):
        \[
        T(x) = \Bigl[\frac{-\eu^{-t^2}}{2t}\Bigr]_x^\infty
        - \int_x^\infty \frac{\eu^{-t^2}}{2t^2}\dd t
        = \frac{\eu^{-x^2}}{2x}
        - \int_x^\infty \frac{\eu^{-t^2}}{2t^2}\dd t .
        \]
        El mismo recurso sobre la nueva integral ($u =
        \frac{-1}{4t^3}$, $v' = -2t\,\eu^{-t^2}$):
        \[
        \int_x^\infty \frac{\eu^{-t^2}}{2t^2}\dd t
        = \frac{\eu^{-x^2}}{4x^3}
        - \frac34\int_x^\infty \frac{\eu^{-t^2}}{t^4}\dd t ,
        \]
        de donde la identidad anunciada.
  \item Una integración por partes más acota el resto:
        \[
        \int_x^\infty \frac{\eu^{-t^2}}{t^4}\dd t
        = \frac{\eu^{-x^2}}{2x^5}
        - \frac52\int_x^\infty\frac{\eu^{-t^2}}{t^6}\dd t
        \leq \frac{\eu^{-x^2}}{2x^5},
        \]
        luego $0 \leq \frac34\int_x^\infty t^{-4}\eu^{-t^2}\dd t
        \leq \frac{3}{8x^5}\eu^{-x^2}$. Descartar el resto
        (positivo) en la identidad de la pregunta 1 da la cota
        inferior; descartar el segundo término (negativo) de la
        primera integración por partes da $T(x) \leq
        \frac{\eu^{-x^2}}{2x}$. Dividiendo el encaje por
        $\frac{\eu^{-x^2}}{2x}$: el cociente queda comprimido entre
        $1 - \frac{1}{2x^2}$ y $1$, luego $T(x) \sim
        \frac{\eu^{-x^2}}{2x}$.
  \item Iterar las integraciones por partes produce la serie
        formal
        \[
        T(x) \approx \frac{\eu^{-x^2}}{2x}\Bigl(1 -
        \frac{1}{2x^2} + \frac{1\cdot3}{(2x^2)^2} -
        \frac{1\cdot3\cdot5}{(2x^2)^3} + \cdots\Bigr),
        \]
        cuyo coeficiente $k$-ésimo $1\cdot3\cdots(2k-1) =
        \frac{(2k)!}{2^k k!}$ crece más deprisa que cualquier
        sucesión geométrica: para $x$ fijo, los términos
        $\frac{1\cdot3\cdots(2k-1)}{(2x^2)^k}$ tienden a infinito
        (su cociente es $\frac{2k+1}{2x^2} \to \infty$), así que la
        serie diverge para todo $x$. Es un desarrollo
        \emph{asintótico}: truncado en cualquier orden fijo, el
        error es del orden del primer término omitido cuando $x \to
        \infty$, pero nunca una serie convergente. (Esta estimación
        de la cola es la cota gaussiana estándar de los capítulos de
        probabilidad.)
\end{enumerate}
\end{solution}

\begin{solution}{pb:b2:integration:1}
\textbf{1.} En $0^+$ el integrando es $\sim t^{x-1}$: la escala de
extremo finito \omterm{def:b2:integration:improper}{converge} si y solo si $1 - x < 1$, es decir, $x >
0$ (y para $x \leq 0$, $t^{x-1} \geq t^{-1}$ diverge); en
$+\infty$, $t^{x-1}\eu^{-t} = o(t^{-2})$ \omterm{def:b2:integration:improper}{converge} para todo $x$.
Después, $\Gamma(x) = \frac{\Gamma(x+1)}{x}$ y $\Gamma(x+1) \to
\Gamma(1) = 1$ cuando $x \to 0^+$ (\omterm{def:b2:metric:continuity}{continuidad},
\cref{thm:b2:integration:gammaprops}): $\Gamma(x) \sim \frac1x$.

\textbf{2.} Con $t = u^2$, $\dd t = 2u\,\dd u$:
\[
\Gamma\Bigl(\frac12\Bigr) = \int_0^\infty
t^{-1/2}\eu^{-t}\dd t
= \int_0^\infty \frac{\eu^{-u^2}}{u}\,2u\,\dd u
= 2\int_0^\infty \eu^{-u^2}\dd u = \sqrt\pi
\]
por el~\cref{exo:b2:integration:8}. Con $u = v/\sqrt2$:
\[
\int_\R \eu^{-v^2/2}\dd v
= 2\sqrt2\int_0^\infty \eu^{-u^2}\dd u
= \sqrt2\,\sqrt\pi = \sqrt{2\pi} .
\]

\textbf{3.} Cierto para $n = 0$ (ambos miembros valen $\sqrt\pi$).
Si $\Gamma(n + \frac12) = \frac{(2n)!}{4^n n!}\sqrt\pi$, la
ecuación funcional da
\[
\Gamma\Bigl(n + 1 + \frac12\Bigr)
= \Bigl(n + \frac12\Bigr)\Gamma\Bigl(n + \frac12\Bigr)
= \frac{2n+1}{2}\cdot\frac{(2n)!}{4^n n!}\sqrt\pi
= \frac{(2n+2)!}{4^{n+1}(n+1)!}\sqrt\pi ,
\]
el último paso porque $\frac{(2n+2)!}{(2n)!} = (2n+2)(2n+1)$ y
$\frac{2n+1}{2} = \frac{(2n+2)(2n+1)}{4(n+1)}$.

\textbf{4.} El~\cref{thm:b2:integration:gammaprops} da $\Gamma''(x)
= \int_0^\infty t^{x-1}\eu^{-t}(\ln t)^2\dd t$ (dos aplicaciones de
la regla de Leibniz, con dominaciones como en la demostración del
teorema); el integrando es $\geq 0$ y no idénticamente nulo, luego
$\Gamma'' > 0$: $\Gamma$ es estrictamente convexa y $\Gamma'$ es
estrictamente creciente. Como $\Gamma(1) = \Gamma(2) = 1$, Rolle
proporciona un $x_0 \in \intoo12$ con $\Gamma'(x_0) = 0$; la
monotonía estricta de $\Gamma'$ hace de $x_0$ su único cero, con
$\Gamma' < 0$ antes y $\Gamma' > 0$ después: $\Gamma$ decrece sobre
$\intoo0{x_0}$, crece sobre $\intoo{x_0}\infty$ y $x_0$ es el
mínimo único.

\textbf{5.} Sea $k \in \N$ y $x \geq 3$; elijamos el entero $n$ con
$n + 1 \leq x < n + 2$ (de modo que $n \geq 1$). Por la monotonía de
la pregunta 4 (válida a partir de $x_0 < 2$): $\Gamma(x) \geq
\Gamma(n + 1) = n!$, mientras que $x^k \leq (n+2)^k$. Como
$\frac{n!}{(n+2)^k} \to \infty$ (los factoriales ganan a las
potencias, volumen del primer año), $\frac{\Gamma(x)}{x^k} \geq
\frac{n!}{(n+2)^k} \to \infty$ cuando $x \to \infty$: $x^k =
o(\Gamma(x))$.

\textbf{6.} Cerca de $0$ el integrando es $\sim t^{x-1}$
(convergente si y solo si $x > 0$), y cerca de $1$ es $\sim
(1-t)^{y-1}$ (si y solo si $y > 0$); ambas comparaciones son entre
funciones positivas, de modo que $B(x,y)$ \omterm{def:b2:integration:improper}{converge} exactamente para
$x, y > 0$. La sustitución $t \mapsto 1 - t$ intercambia los dos
factores: $B(x,y) = B(y,x)$.

\textbf{7.} $B(x,1) = \int_0^1 t^{x-1}\dd t = \frac1x$. Por partes
sobre $\intcc\varepsilon{1-\varepsilon}$ con $u = (1-t)^y$ y $v =
\frac{t^x}{x}$:
\[
\int t^{x-1}(1-t)^{y}\dd t
= \Bigl[\frac{t^x(1-t)^y}{x}\Bigr]
+ \frac{y}{x}\int t^{x}(1-t)^{y-1}\dd t ;
\]
el corchete se anula en ambos extremos cuando $\varepsilon \to 0$
($x > 0$ en $0$, $y > 0$ en $1$), dejando $B(x, y+1) = \frac
yx\,B(x+1, y)$.

\textbf{8.} Como $t + (1-t) = 1$:
\[
t^{x-1}(1-t)^{y-1} = t^{x}(1-t)^{y-1} + t^{x-1}(1-t)^{y},
\]
luego $B(x,y) = B(x+1,y) + B(x,y+1)$. La pregunta 7 se lee
$B(x+1,y) = \frac xy B(x,y+1)$; sustituyendo,
\[
B(x,y) = \Bigl(\frac xy + 1\Bigr)B(x,y+1)
= \frac{x+y}{y}\,B(x,y+1),
\]
es decir, $B(x,y+1) = \frac{y}{x+y}B(x,y)$; la relación gemela se
sigue de la simetría de la pregunta 6.

\textbf{9.} Inducción sobre $n$ con $m$ fijo: $B(m,1) = \frac1m =
\frac{(m-1)!\,0!}{m!}$ y, si la fórmula vale en $n$,
\[
B(m, n+1) = \frac{n}{m+n}\,B(m,n)
= \frac{n}{m+n}\cdot\frac{(m-1)!(n-1)!}{(m+n-1)!}
= \frac{(m-1)!\,n!}{(m+n)!} .
\]
Reescribiendo: $B(m,n) = \frac{(m-1)!(n-1)!}{(m+n-1)!} =
\bigl[(m+n-1)\binom{m+n-2}{m-1}\bigr]^{-1}$.

\textbf{10.} Ambos miembros de $B(x,n) =
\frac{\Gamma(x)\Gamma(n)}{\Gamma(x+n)}$ valen $\frac1x$ en $n = 1$
($\Gamma(1) = 1$, $\Gamma(x+1) = x\Gamma(x)$). Si coinciden en $n$,
entonces, por la relación de descenso y la ecuación funcional:
\[
B(x, n+1) = \frac{n}{x+n}\,B(x,n), \qquad
\frac{\Gamma(x)\Gamma(n+1)}{\Gamma(x+n+1)}
= \frac{n}{x+n}\cdot
\frac{\Gamma(x)\Gamma(n)}{\Gamma(x+n)} :
\]
las dos sucesiones obedecen la misma recursión desde la misma
semilla, luego coinciden para todo $n \in \N^*$ y todo $x > 0$.

\textbf{11.} Con $t = \sin^2\theta$ ($\theta \in \intoo0{\pi/2}$,
$\dd t = 2\sin\theta\cos\theta\,\dd\theta$), $t^{x-1} =
\sin^{2x-2}\theta$ y $(1-t)^{y-1} = \cos^{2y-2}\theta$:
\[
B(x,y) = \int_0^{\pi/2}\sin^{2x-2}\theta\,\cos^{2y-2}\theta
\cdot 2\sin\theta\cos\theta\,\dd\theta
= 2\int_0^{\pi/2}\sin^{2x-1}\theta\,\cos^{2y-1}\theta\,
\dd\theta .
\]

\textbf{12.} Tomemos $y = \frac12$ (lo que mata el factor coseno) y
$2x - 1 = n$: $B\bigl(\frac{n+1}2, \frac12\bigr) = 2W_n$, es decir,
$W_n = \frac12 B\bigl(\frac{n+1}2,\frac12\bigr)$. La relación de
descenso en la primera variable da
\[
\frac{W_n}{W_{n-2}}
= \frac{B\bigl(\frac{n-1}2 + 1, \frac12\bigr)}
{B\bigl(\frac{n-1}2, \frac12\bigr)}
= \frac{\frac{n-1}2}{\frac{n-1}2 + \frac12}
= \frac{n-1}{n} :
\]
la recurrencia de Wallis, esta vez sin ninguna integración por
partes sobre senos; la parte II hizo el trabajo de una vez por
todas.

\textbf{13.} $B\bigl(\frac12,\frac12\bigr) = 2W_0 =
2\cdot\frac\pi2 = \pi$, mientras que
$\Gamma\bigl(\frac12\bigr)^2/\Gamma(1) = (\sqrt\pi)^2 = \pi$: la
fórmula de Euler se cumple en $\bigl(\frac12,\frac12\bigr)$.

\textbf{14.} Iterando $W_{2n} = \frac{2n-1}{2n}W_{2n-2}$ desde $W_0
= \frac\pi2$:
\[
W_{2n} = \frac\pi2\prod_{k=1}^{n}\frac{2k-1}{2k}
= \frac\pi2\cdot\frac{(2n)!}{4^n(n!)^2},
\]
ya que $\prod(2k-1) = \frac{(2n)!}{2^n n!}$ y $\prod 2k = 2^n n!$.
Por tanto, usando la pregunta 3:
\[
B\Bigl(n+\frac12, \frac12\Bigr) = 2W_{2n}
= \pi\,\frac{(2n)!}{4^n(n!)^2}
= \frac{(2n)!\sqrt\pi}{4^n n!}\cdot\frac{\sqrt\pi}{n!}
= \frac{\Gamma\bigl(n+\frac12\bigr)\Gamma\bigl(\frac12\bigr)}
{\Gamma(n+1)} .
\]
Fijemos ahora $x \in \frac12\N^*$. La fórmula de Euler vale en $(x,
\frac12)$: para $x$ entero es la pregunta 10 (con la simetría), y
para $x = n + \frac12$ es la fórmula anterior. Ambos miembros de la
fórmula de Euler obedecen la recursión de descenso $y \mapsto y + 1$
(la pregunta 8 a la izquierda y la ecuación funcional a la derecha,
como en la pregunta 10): la inducción propaga la fórmula desde $y =
\frac12$ e $y = 1$ hasta todo $y \in \frac12\N^*$. La fórmula de
Euler vale, pues, siempre que $2x, 2y \in \N^*$.

\textbf{15.} Con $u = \frac{t}{1-t}$, es decir, $t =
\frac{u}{1+u}$, $1 - t = \frac{1}{1+u}$ y $\dd t = \frac{\dd
u}{(1+u)^2}$:
\[
B(x,y) = \int_0^\infty
\Bigl(\frac{u}{1+u}\Bigr)^{x-1}
\Bigl(\frac{1}{1+u}\Bigr)^{y-1}
\frac{\dd u}{(1+u)^2}
= \int_0^\infty \frac{u^{x-1}}{(1+u)^{x+y}}\,\dd u .
\]
En $x = y = \frac12$, con $u = v^2$:
\[
\int_0^\infty \frac{u^{-1/2}}{1+u}\dd u
= \int_0^\infty \frac{2\,\dd v}{1+v^2} = \pi
= B\Bigl(\frac12,\frac12\Bigr) . \checkmark
\]

\textbf{16.} Una integración por partes, para $1 \leq k \leq n$ y
$s > 0$ ($u = (1 - t/n)^k$, $v = t^s/s$; los términos de frontera se
anulan):
\[
\int_0^n \Bigl(1-\frac tn\Bigr)^{\!k} t^{s-1}\dd t
= \frac{k}{ns}\int_0^n \Bigl(1-\frac tn\Bigr)^{\!k-1}
t^{s}\dd t .
\]
Partiendo de $k = n$, $s = x$ e iterando $n$ veces:
\[
\int_0^n \Bigl(1-\frac tn\Bigr)^{\!n} t^{x-1}\dd t
= \frac{n(n-1)\cdots1}{n^n\,x(x+1)\cdots(x+n-1)}
\int_0^n t^{x+n-1}\dd t
= \frac{n!}{n^n}\cdot
\frac{n^{x+n}}{x(x+1)\cdots(x+n)} ,
\]
que vale $\dfrac{n!\,n^x}{x(x+1)\cdots(x+n)}$.

\textbf{17.} Por el~\cref{exo:b2:integration:7}, el miembro
izquierdo tiende a $\Gamma(x)$ (convergencia dominada con dominador
$t^{x-1}\eu^{-t}$); y el derecho es el cociente de Gauss:
\[
\Gamma(x) = \lim_{n\to\infty}
\frac{n!\,n^x}{x(x+1)\cdots(x+n)} .
\]

\textbf{18.} Tomando logaritmos en el cociente $G_n(x)$ de la
pregunta 16 y separando $\ln(x+k) = \ln k + \ln(1 + x/k)$ para $k
\geq 1$:
\[
\ln G_n(x) = x\ln n - \ln x - \sum_{k=1}^n
\ln\Bigl(1+\frac xk\Bigr)
= -\ln x + x(\ln n - H_n)
+ \sum_{k=1}^n\Bigl(\frac xk -
\ln\Bigl(1+\frac xk\Bigr)\Bigr).
\]
Para $u \geq 0$ se tiene $u - \frac{u^2}2 \leq \ln(1+u) \leq u$, de
modo que el término general está en $\intcc{0}{x^2/(2k^2)}$: la
serie \omterm{def:b2:integration:improper}{converge} (comparación con $\sum k^{-2}$). Como $\ln n - H_n
\to -\gamma$ (\cref{ex:b2:comparison:harmonic}) y $\ln G_n(x) \to
\ln\Gamma(x)$ (pregunta 17 y \omterm{def:b2:metric:continuity}{continuidad} de $\ln$):
\[
\ln\Gamma(x) = -\ln x - \gamma x
+ \sum_{k=1}^\infty\Bigl(\frac xk -
\ln\Bigl(1+\frac xk\Bigr)\Bigr) .
\]

\textbf{19.} En $x = \frac12$, el denominador es
$\prod_{k=0}^n\bigl(k+\frac12\bigr) =
\frac{(2n+1)!}{2^{2n+1}n!}$ (desarrollando las mitades), luego
\[
G_n\Bigl(\frac12\Bigr)
= \frac{n!\,\sqrt n\;2^{2n+1}n!}{(2n+1)!}
= \frac{2\sqrt n\;4^n}{(2n+1)\binom{2n}{n}} .
\]
Con $\binom{2n}{n} \sim \frac{4^n}{\sqrt{\pi n}}$
(\cref{ex:b2:comparison:centralbinomial}):
\[
G_n\Bigl(\frac12\Bigr)
\sim \frac{2\sqrt n\,\sqrt{\pi n}}{2n+1}
\longrightarrow \sqrt\pi = \Gamma\Bigl(\frac12\Bigr) .
\]
El $\sqrt\pi$ del coeficiente binomial central (que venía de
Wallis y, por tanto, de la constante de Stirling) y el $\sqrt\pi$
de la integral de Gauss son el mismo número.

\textbf{20.} La identidad es álgebra directa: multiplíquese
$\frac{n!\,n^x}{x(x+1)\cdots(x+n)}$ por $\frac{nx}{x+n+1}$ y
absórbase $x$ en el producto y $n$ en $n^x$. Haciendo $n \to
\infty$: el miembro izquierdo tiende a $\Gamma(x+1)$ (Gauss en
$x+1$) y el derecho a $\Gamma(x)\cdot x\cdot 1$, ya que
$\frac{n}{x+n+1} \to 1$: $\Gamma(x+1) = x\Gamma(x)$, recuperado
sin una sola integración por partes.

\textbf{21.} Con $u = t^a$, $t = u^{1/a}$, $\dd t = \frac1a u^{1/a
- 1}\dd u$:
\[
\int_0^\infty \eu^{-t^a}\dd t
= \frac1a\int_0^\infty u^{\frac1a - 1}\eu^{-u}\dd u
= \frac1a\,\Gamma\Bigl(\frac1a\Bigr)
= \Gamma\Bigl(1 + \frac1a\Bigr) .
\]
Cuando $a \to +\infty$ (a lo largo de cualquier sucesión):
$\eu^{-t^a} \to 1$ para $0 < t < 1$, $\to \eu^{-1}$ en $t = 1$ y
$\to 0$ para $t > 1$; para $a \geq 2$ domínese por $\mathbf 1_{t
\leq 1} + \eu^{-t^2}\mathbf 1_{t > 1}$ ($t^a \geq t^2$ para $t
\geq 1$), integrable. Por convergencia dominada, la integral
tiende a $\int_0^1 1\,\dd t = 1$, como debe ser, ya que $\Gamma(1
+ \frac1a) \to \Gamma(1) = 1$ por \omterm{def:b2:metric:continuity}{continuidad}.

\textbf{22.} Con $u = t^n$, $\dd t = \frac1n u^{1/n - 1}\dd u$:
\[
\int_0^1 \frac{\dd t}{\sqrt{1-t^n}}
= \frac1n\int_0^1 u^{\frac1n-1}(1-u)^{-1/2}\dd u
= \frac1n\,B\Bigl(\frac1n, \frac12\Bigr) .
\]
$n = 1$: $B\bigl(1,\frac12\bigr) = B\bigl(\frac12,1\bigr) = 2$, en
concordancia con $\int_0^1\frac{\dd t}{\sqrt{1-t}} = 2$. $n = 2$:
$\frac12 B\bigl(\frac12,\frac12\bigr) = \frac\pi2 = \arcsin 1$.
Para $n = 4$, el valor $\frac14 B\bigl(\frac14,\frac12\bigr)$ es la
constante de la lemniscata: sin forma cerrada elemental.

\textbf{23.} Iterando la ecuación funcional:
\[
\frac{1}{\Gamma(x)}\int_0^\infty t^{x+k-1}\eu^{-t}\dd t
= \frac{\Gamma(x+k)}{\Gamma(x)}
= (x+k-1)(x+k-2)\cdots x ,
\]
el factorial ascendente con $k$ factores. En $x = 1$:
$\Gamma(1+k)/\Gamma(1) = k!$, los momentos de $\eu^{-t}$ del~\cref{exo:b2:integration:2}.

\textbf{24.} Sustituyamos $t = \frac{1+s}2$ ($s \in \intoo{-1}1$,
$\dd t = \frac{\dd s}2$, $t(1-t) = \frac{1-s^2}4$):
\[
B(x,x) = \int_{-1}^{1}\Bigl(\frac{1-s^2}{4}\Bigr)^{x-1}
\frac{\dd s}{2}
= 4^{1-x}\int_0^1 (1-s^2)^{x-1}\dd s
\]
(el integrando es par). Después, $s = \sqrt v$ ($\dd s = \frac{\dd
v}{2\sqrt v}$):
\[
B(x,x) = \frac{4^{1-x}}{2}\int_0^1
v^{-1/2}(1-v)^{x-1}\dd v
= 2^{1-2x}\,B\Bigl(\frac12, x\Bigr) .
\]
Para $2x \in \N^*$, todos los argumentos a la vista están en
$\frac12\N^*$, así que la fórmula de Euler (pregunta 14) se aplica
a ambos miembros:
\[
\frac{\Gamma(x)^2}{\Gamma(2x)}
= 2^{1-2x}\,
\frac{\Gamma\bigl(\frac12\bigr)\Gamma(x)}
{\Gamma\bigl(x+\frac12\bigr)}
\quad\Longleftrightarrow\quad
\Gamma(x)\,\Gamma\Bigl(x+\frac12\Bigr)
= 2^{1-2x}\sqrt\pi\;\Gamma(2x) .
\]
Comprobación directa en $x = n$: el miembro izquierdo vale
$(n-1)!\cdot \frac{(2n)!\sqrt\pi}{4^n n!} = \frac{(2n)!\sqrt\pi}{4^n
n}$, y el derecho $2\cdot4^{-n}\sqrt\pi\,(2n-1)! =
\frac{(2n)!\sqrt\pi}{4^n n}$: iguales.

\textbf{25.} (i) La integración por partes produjo $B(x,y+1) =
\frac yx B(x+1,y)$, la única identidad de la que manan todas las
relaciones de descenso, los valores enteros y semienteros y la
recurrencia de Wallis. (ii) La convergencia dominada convirtió las
integrales elementales $\int_0^n(1-t/n)^n t^{x-1}$ en $\Gamma(x)$
(fórmula de Gauss, pregunta 17) y calculó el límite $a \to \infty$
de la pregunta 21. (iii) Del capítulo de comparación importamos la
constante de Euler ($\ln n - H_n \to -\gamma$, pregunta 18) y la
asintótica del binomial central (pregunta 19), es decir, la fórmula
de Stirling disfrazada. (iv) La fórmula de Euler $B(x,y) =
\Gamma(x)\Gamma(y)/\Gamma(x+y)$ queda demostrada para $y \in \N^*$
con $x > 0$ arbitrario (pregunta 10) y para todos los pares
semienteros (pregunta 14); el caso general $x, y > 0$ espera al
cálculo por integral doble del capítulo de integrales múltiples,
que factoriza $\Gamma(x)\Gamma(y)$ sobre un cuadrante.
\end{solution}
